\section{Introduction}

An index does not release space when rows are deleted from the table it
indexes. Reclaiming that space is the job of a separate maintenance process,
which must find the entries pointing at deleted rows, remove them, notice pages
that have become empty and return those pages to use~--- concurrently with the
ordinary workload, without blocking it.

For a B-tree this is well-trodden ground~\cite{graefe2011modern}. For a
generalized search tree~\cite{hellerstein1995generalized} it is harder in three
specific ways, and each of them was, until recently, unsolved in the most
widely deployed implementation. There is no order on keys, so the structure
offers no natural sequence in which to walk it. The descent is multi-path, so
the parent of a page is not known when the page is reached other than by
descending to it. And the key of a parent cannot be recomputed locally, so
removing an entry is not a local operation.

This paper reports the algorithms that address these points, their
implementation in PostgreSQL, and five years of their operation. It also
reports, at some length, a case in which an algorithm from this same line of
work was published, reviewed, committed, and then reverted after a bug report
from production. That case is not an aside: it is the reason the paper ends
with a claim about what peer review of a concurrent algorithm can and cannot
establish.

\paragraph{Contributions.}

\begin{enumerate}
\item Garbage collection that traverses index pages in physical rather than
      logical order, and the treatment of entries that concurrent insertions
      move behind the sweep (Section~\ref{sec:physical}).
\item A two-pass algorithm that deletes emptied pages and returns them to the
      free space map, with the lock coupling that makes it safe against
      concurrent descent (Section~\ref{sec:deletion}).
\item The safety condition for reusing a deleted page, expressed over full
      64-bit transaction identifiers rather than wrapping 32-bit ones, and the
      failure mode that the wrapping version exhibits~--- not corruption, but
      the silent cessation of page reuse after long operation
      (Section~\ref{sec:xid}).
\item An account of the reverted lock-granularity change: what was published,
      what production reported, why no non-invasive fix was found, and what
      follows for how such work should be validated (Section~\ref{sec:revert}).
\end{enumerate}

\section{Background}
\label{sec:background}

% Обобщённое дерево поиска, схема Корнакера, правые ссылки и номера
% последовательности узлов. Многоверсионность: удалённая запись остаётся, пока
% её может увидеть открытое сканирование. Понятие «страница больше не видна ни
% одному сканированию» и почему оно нетривиально при многопутевом спуске.

\section{Traversal in physical order}
\label{sec:physical}

% Логический обход даёт произвольное чтение: P = (L+I) c_r.
% Физический: P = (L+I) c_s + beta c_r, где beta — страницы, перечитанные из-за
% конкурентных вставок, переместивших необработанные элементы на уже пройденные
% страницы.
% Затруднение: при физическом обходе неизвестен родитель, а он нужен, чтобы
% удалить нисходящую ссылку. Отсюда два прохода (раздел \ref{sec:deletion}).
% Коммит fe280694d0d, выпуск 12.

\section{Deleting and reusing pages}
\label{sec:deletion}

% До этой работы обобщённое дерево поиска не переиспользовало страницы вовсе:
% опустевшая страница оставалась навсегда. При нагрузке «вставка в новую
% область, удаление из старой» индекс рос неограниченно.
%
% Двухпроходный алгоритм: (1) физический обход, сбор номеров опустевших листьев
% и внутренних страниц; (2) обход внутренних, удаление нисходящих ссылок,
% пометка страниц удалёнными.
%
% Порядок операций внутри второго прохода — условие корректности: сначала
% удаляется ссылка из родителя, затем страница помечается удалённой, блокировки
% сцепляются. Обратный порядок допускал бы вставку в удаляемую страницу.
%
% Ограничение: внутренние страницы не удаляются, последний потомок не удаляется
% никогда, как и в B-дереве. Рост замедляется на порядок, но не прекращается.
% Коммиты 7df159a620b, 9eb5607e699, выпуск 12.

\section{When is it safe to reuse a page}
\label{sec:xid}

% Сканирование, начатое до удаления, может обратиться к странице позже.
% Условие: «страница удалена раньше, чем начались все ныне активные
% транзакции».
%
% При 32-разрядном счётчике условие теряет смысл при зацикливании: достаточно
% старая страница снова выглядит недавно удалённой. Отказ не наступает — просто
% прекращается возврат страниц в оборот, и обнаруживается это спустя месяцы.
%
% Общая формулировка: условия безопасности, выражаемые через монотонный
% счётчик, должны формулироваться в терминах счётчика, не подверженного
% зацикливанию, иначе корректность зависит от длительности эксплуатации.
% Коммит 6655a7299d8, выпуск 13, перенесено в 12.
% Смежное: 4ed8f0913bf, 64-разрядная адресация служебных журналов.

\section{Lock granularity, and a result that did not survive}
\label{sec:revert}

% Ключевой раздел. Изложение по порядку:
%
% 1. Задача: сборка мусора в дереве вхождений обратного индекса захватывала
%    исключительную блокировку корня на всё время обработки. Для частого ключа
%    дерево велико, и блокировка останавливает все обращения к этому ключу.
%
% 2. Предложено два изменения: (а) захватывать блокировку только когда
%    действительно есть что удалять; (б) блокировать поддерево, а не всё дерево.
%    Коммит 218f51584d5, выпуск 10. Опубликовано [borodin2018gin].
%
% 3. Через полтора года: сообщение об ошибке из эксплуатации
%    [bug-gin-lwlock-2018] — соединения зависают на LWLock buffer_content.
%    Диагноз: при конкурентном расщеплении родителя вставка не удерживает
%    закреплений всех страниц пути корень—лист, откуда цикл ожидания со
%    сборкой мусора.
%
% 4. Неинвазивного решения не нашли. Часть (б) отменена коммитом fd83c83d094 с
%    переносом во все поддерживаемые выпуски. Часть (а) сохранена и даёт
%    основной практический эффект: в большинстве проходов удалять нечего.
%
% 5. Сопутствующее: 52ac6cd2d0c — с версии 9.4 сканирование спускается не от
%    корня, поэтому приём «переход вправо держит две страницы» перестал быть
%    достаточным, и страница могла быть переиспользована до обращения к ней.
%    Момент удаления записан в pd_prune_xid, поскольку формат служебной области
%    страницы менять нельзя из соображений двоичной совместимости.
%
% 6. Вывод, ради которого раздел и написан: рецензирование статьи проверяет
%    рассуждение при заявленных предпосылках. Предпосылки конкурентного
%    алгоритма относятся не к нему, а к остальной части системы — здесь к
%    протоколу закреплений при вставке, — и их нарушение обнаруживается только
%    эксплуатацией на разнообразных нагрузках. Отсюда два практических
%    следствия: (1) публикация о конкурентном алгоритме должна называть
%    предпосылки в проверяемой форме; (2) свойства, не проверяемые
%    функциональным тестом, должны проверяться инвариантами структуры.

\section{Prefetching during maintenance}
\label{sec:prefetch}

% Выигрыш физического порядка пропорционален разности стоимостей произвольного
% и последовательного чтения. На современных накопителях эта разность
% сокращается, но появляется другой источник: зная последовательность страниц
% заранее, можно выдавать несколько запросов чтения одновременно.
% Коммиты c5c239e26e3, 69273b818b1, e215166c9c8, выпуск 18.
% Страницы, требующие повторной обработки (beta), читаются обычным образом:
% их последовательность заранее неизвестна.

\section{Evaluation}
\label{sec:eval}

\subsection{Traversal order}

The pair of builds differs by exactly one commit. Ten million points, every
fifth row deleted, index of about 90\,800 pages, buffer cache deliberately set
to 256~MB so that the index does not fit: with a cache that holds the index,
traversal order cannot matter. Each measurement starts from a fresh copy of an
unchanged master table.

\begin{table}[t]
\centering
\caption{Logical against physical traversal order}
\label{tab:traversal}
\begin{tabular}{lrrr}
\toprule
traversal & VACUUM, s & index pages read & index pages \\
\midrule
logical & 12.1 / 11.8 / 12.6 & 180\,020 / 179\,772 / 180\,028 & 90\,827 \\
physical & 11.8 / 11.7 & 90\,045 / 90\,013 & 90\,870 \\
\bottomrule
\end{tabular}
\end{table}

Physical order halves the number of index page reads, exactly and
reproducibly: depth-first traversal returns to pages it has already read, and
with a cache of 32\,768 pages against an index of 90\,800 those pages are gone
by the time it does. The ratio is 2.00 in every run.

End-to-end time does not change. The reason is that this workload is dominated
by the heap scan, which both versions perform identically; the saving falls
entirely on the index half of the work. We report it as such rather than as a
speedup of VACUUM, which it is not: whether it shows up in wall-clock time
depends on what fraction of the cleanup is index work.

This is also why the primary quantity in this paper is the number of index page
reads and not elapsed time. The former isolates the change; the latter mixes it
with work the change does not touch.

% TODO: остальные измерения
% 2. Возврат страниц: 7df159a620b^ против 7df159a620b на нагрузке «вставка в
%    новую область, удаление из старой». Динамика объёма индекса во времени.
% 3. Упреждающее чтение: выпуск 17 против 18 за пределами буферного кэша.
% 4. Блокировки обратного индекса: время ожидания, 218f51584d5^ против
%    218f51584d5 и против fd83c83d094.


\section{Deployment experience}
\label{sec:deployment}

% Пять лет в выпусках 12–19. Что подтвердилось эксплуатацией и что она выявила:
% три исправления взаимных блокировок после первой редакции изменения
% гранулярности (fd83c83d094, c6ade7a8cd3, e14641197a5) — все обнаружены
% сообщениями об ошибках, ни одно тестом.
%
% Это и есть содержание раздела: перечислить, что было найдено эксплуатацией и
% чего не нашло бы тестирование.

\section{Limitations}

% Возврат страниц не полон: внутренние страницы и последний потомок не
% удаляются. Выигрыш физического обхода зависит от носителя. Двухпроходный
% алгоритм требует памяти под карту опустевших страниц; при её нехватке нужен
% запасной путь. Блокировка поддерева при действующем протоколе закреплений
% некорректна и требует его пересмотра — задача открыта.

\section{Conclusion}
